\documentclass[10pt,a4paper]{amsart}
\usepackage{amsmath,amssymb,amsthm,graphicx}
\usepackage{verbatim}
\newenvironment{reference}{\comment}{\endcomment}
%\newenvironment{reference}{}{}
\newenvironment{slogan}{\comment}{\endcomment}
\newenvironment{history}{\comment}{\endcomment}

% For commutative diagrams we use Xy-pic
\usepackage[all]{xy}

% We use 2cell for 2-commutative diagrams.
\xyoption{2cell}
% Theorem environments.
%
\theoremstyle{plain}
\newtheorem{theorem}[subsection]{Theorem}
\newtheorem{proposition}[subsection]{Proposition}
\newtheorem{lemma}[subsection]{Lemma}

\theoremstyle{definition}
\newtheorem{definition}[subsection]{Definition}
\newtheorem{example}[subsection]{Example}
\newtheorem{exercise}[subsection]{Exercise}
\newtheorem{situation}[subsection]{Situation}

\theoremstyle{remark}
\newtheorem{remark}[subsection]{Remark}
\newtheorem{remarks}[subsection]{Remarks}

\numberwithin{equation}{subsection}

% Macros
%
\def\lim{\mathop{\mathrm{lim}}\nolimits}
\def\colim{\mathop{\mathrm{colim}}\nolimits}
\def\Spec{\mathop{\mathrm{Spec}}}
\def\Hom{\mathop{\mathrm{Hom}}\nolimits}
\def\Ext{\mathop{\mathrm{Ext}}\nolimits}
\def\SheafHom{\mathop{\mathcal{H}\!\mathit{om}}\nolimits}
\def\SheafExt{\mathop{\mathcal{E}\!\mathit{xt}}\nolimits}
\def\Sch{\mathit{Sch}}
\def\Mor{\mathop{\mathrm{Mor}}\nolimits}
\def\Ob{\mathop{\mathrm{Ob}}\nolimits}
\def\Sh{\mathop{\mathit{Sh}}\nolimits}
\def\NL{\mathop{N\!L}\nolimits}
\def\CH{\mathop{\mathrm{CH}}\nolimits}
\def\proetale{{pro\text{-}\acute{e}tale}}
\def\etale{{\acute{e}tale}}
\def\QCoh{\mathit{QCoh}}
\def\Ker{\mathop{\mathrm{Ker}}}
\def\Im{\mathop{\mathrm{Im}}}
\def\Coker{\mathop{\mathrm{Coker}}}
\def\Coim{\mathop{\mathrm{Coim}}}

% Boxtimes
%
\DeclareMathSymbol{\boxtimes}{\mathbin}{AMSa}{"02}

%
% Macros for moduli stacks/spaces
%
\def\QCohstack{\mathcal{QC}\!\mathit{oh}}
\def\Cohstack{\mathcal{C}\!\mathit{oh}}
\def\Spacesstack{\mathcal{S}\!\mathit{paces}}
\def\Quotfunctor{\mathrm{Quot}}
\def\Hilbfunctor{\mathrm{Hilb}}
\def\Curvesstack{\mathcal{C}\!\mathit{urves}}
\def\Polarizedstack{\mathcal{P}\!\mathit{olarized}}
\def\Complexesstack{\mathcal{C}\!\mathit{omplexes}}
% \Pic is the operator that assigns to X its picard group, usage \Pic(X)
% \Picardstack_{X/B} denotes the Picard stack of X over B
% \Picardfunctor_{X/B} denotes the Picard functor of X over B
\def\Pic{\mathop{\mathrm{Pic}}\nolimits}
\def\Picardstack{\mathcal{P}\!\mathit{ic}}
\def\Picardfunctor{\mathrm{Pic}}
\def\Deformationcategory{\mathcal{D}\!\mathit{ef}}
\usepackage{cleveref}
\crefname{originalreference}{original reference}{original references}
\Crefname{originalreference}{Original reference}{Original references}
\makeatletter
\newlabel{situation-pre-pure}{{05FW}{source}{}{sourcepointer-1}{}}
\newlabel{situation-pre-pure@cref}{{[originalreference][0][]05FW}{[1][0][]source}{}{}{}}
\makeatother
\crefname{originallemma}{original lemma}{original lemmas}
\Crefname{originallemma}{Original lemma}{Original lemmas}
\makeatletter
\newlabel{spaces-divisors-lemma-base-change-relative-assassin}{{0CV2}{source}{}{sourcepointer-2}{}}
\newlabel{spaces-divisors-lemma-base-change-relative-assassin@cref}{{[originallemma][0][]0CV2}{[1][0][]source}{}{}{}}
\makeatother
\crefname{originallemma}{original lemma}{original lemmas}
\Crefname{originallemma}{Original lemma}{Original lemmas}
\makeatletter
\newlabel{spaces-limits-lemma-inverse-limit-irreducibles}{{0CUK}{source}{}{sourcepointer-3}{}}
\newlabel{spaces-limits-lemma-inverse-limit-irreducibles@cref}{{[originallemma][0][]0CUK}{[1][0][]source}{}{}{}}
\makeatother
\crefname{originallemma}{original lemma}{original lemmas}
\Crefname{originallemma}{Original lemma}{Original lemmas}
\makeatletter
\newlabel{spaces-limits-lemma-limit-nonempty}{{086W}{source}{}{sourcepointer-4}{}}
\newlabel{spaces-limits-lemma-limit-nonempty@cref}{{[originallemma][0][]086W}{[1][0][]source}{}{}{}}
\makeatother
\title{An impurity descends to some stage of a directed limit of affine schemes}
\author{The Stacks Project Authors}
\date{stacks-project:spaces-flat.tex:local-snapshot-20261012; GNU FDL 1.2 or later}
\begin{document}
\maketitle
\noindent\textit{Editorial note.} All delivered mathematical text is the author's own, copied from the pinned Stacks Project chapter file spaces-flat.tex; no statement or proof step was rewritten, reordered or removed.\par\medskip
\noindent\textit{Editorial note.} The counted proof is the author's complete argument: a base change verification of the relative assassin condition followed by a limit and quasi-compactness argument producing a stage where the closure condition already holds\par\medskip
\noindent\textit{Editorial note.} The label lemma-impure-limit names the counted statement, and the proof environment immediately following it in the source is delivered verbatim as its proof.\par\medskip
\noindent\textit{Editorial note.} The counted statement is carried by the author's own lemma environment in the chapter source.\par\medskip
\noindent\textit{Allowed standard prerequisites (editorial).}\begin{itemize}
\item Impurities above a point as witnesses of impurity, and the base change of the relative assassin of a module
\item The reduced induced structure on the closure of a point, and its behaviour under a directed limit of the base
\item Quasi-compactness of the fibres of a morphism of finite type, and the criterion that a directed limit of quasi-compact spaces is nonempty exactly when some stage is
\end{itemize}\medskip
\ifdefined\hypertarget\hypertarget{sourcepointer-1}{}\fi
\noindent\textit{Source reference.} Original reference 05FW: sources/stacks-spaces-flat/pinned-original-source.tex, source lines 260--260. Prerequisite pointer invoked at proof line(s) 320; it names a standard fact of the same chapter that the counted argument uses, and every step of the counted proof built on it is retained verbatim in the delivered body.\par\medskip
\ifdefined\hypertarget\hypertarget{sourcepointer-2}{}\fi
\noindent\textit{Source reference.} Original lemma 0CV2: sources/stacks-chow/tags.txt, source lines 16122--16122. Prerequisite pointer invoked at proof line(s) 333; it names a standard fact of a sibling chapter that the counted argument uses, and every step of the counted proof built on it is retained verbatim in the delivered body.\par\medskip
\ifdefined\hypertarget\hypertarget{sourcepointer-3}{}\fi
\noindent\textit{Source reference.} Original lemma 0CUK: sources/stacks-chow/tags.txt, source lines 16105--16105. Prerequisite pointer invoked at proof line(s) 344; it names a standard fact of a sibling chapter that the counted argument uses, and every step of the counted proof built on it is retained verbatim in the delivered body.\par\medskip
\ifdefined\hypertarget\hypertarget{sourcepointer-4}{}\fi
\noindent\textit{Source reference.} Original lemma 086W: sources/stacks-chow/tags.txt, source lines 10406--10406. Prerequisite pointer invoked at proof line(s) 357; it names a standard fact of a sibling chapter that the counted argument uses, and every step of the counted proof built on it is retained verbatim in the delivered body.\par\medskip
\begin{lemma}
\label{lemma-impure-limit}
In Situation \ref{situation-pre-pure}.
Let $(g : T \to S, t' \leadsto t, \xi)$ be an impurity of
$\mathcal{F}$ above $y$. Assume $T = \lim_{i \in I} T_i$ is a directed limit
of affine schemes over $Y$. Then for some $i$ the triple
$(T_i \to Y, t'_i \leadsto t_i, \xi_i)$ is an impurity of
$\mathcal{F}$ above $y$.
\end{lemma}

\begin{proof}
The notation in the statement means this: Let $p_i : T \to T_i$
be the projection morphisms, let $t_i = p_i(t)$ and $t'_i = p_i(t')$.
Finally $\xi_i \in |X_{T_i}|$ is the image of $\xi$. By
Divisors on Spaces, Lemma
\ref{spaces-divisors-lemma-base-change-relative-assassin}
we have $\xi_i \in \text{Ass}_{X_{T_i}/T_i}(\mathcal{F}_{T_i})$.
Thus the only point is to show that
$t_i \not \in f_{T_i}(\overline{\{\xi_i\}})$ for some $i$.

\medskip\noindent
Let $Z_i \subset X_{T_i}$ be the reduced induced scheme structure
on $\overline{\{\xi_i\}} \subset |X_{T_i}|$
and let $Z \subset X_T$ be the reduced induced scheme structure on
$\overline{\{\xi\}} \subset |X_T|$.
Then $Z = \lim Z_i$ by
Limits of Spaces, Lemma \ref{spaces-limits-lemma-inverse-limit-irreducibles}
(the lemma applies because each $X_{T_i}$ is decent).
Choose a field $k$ and a morphism $\Spec(k) \to T$ whose image is $t$.
Then
$$
\emptyset =
Z \times_T \Spec(k) = (\lim Z_i) \times_{(\lim T_i)} \Spec(k)
= \lim Z_i \times_{T_i} \Spec(k)
$$
because limits commute with fibred products (limits commute with limits).
Each $Z_i \times_{T_i} \Spec(k)$ is quasi-compact because $X_{T_i} \to T_i$
is of finite type and hence $Z_i \to T_i$ is of finite type.
Hence $Z_i \times_{T_i} \Spec(k)$ is empty for some $i$ by
Limits of Spaces, Lemma \ref{spaces-limits-lemma-limit-nonempty}.
Since the image of the composition $\Spec(k) \to T \to T_i$ is $t_i$
we obtain what we want.
\end{proof}

\end{document}
